The central claim of the Codynamic Book Machine is that a book is not produced by a single pass from outline to PDF, but by a coupled system in which structure, prose, agents, queues, dependencies, artifacts, and compiled outputs remain in motion together. In this view, the manuscript is not a static target that agents merely fill in; it is a codynamic system whose parts continuously constrain and update one another.

The outline gives the system its intended shape, but the outline is only one state surface among several. Section drafts, task queues, runtime messages, registered dependencies, reference records, diagram assets, and build outputs all participate in the same evolving work. When one of these surfaces changes, the others may need to change as well. A revised section can expose a missing citation. A new dependency can alter the order in which sections should be drafted. A build failure can reveal a structural problem in the TeX payload. A gardener pass can identify drift between a section and its siblings. The machine is designed so that these changes are not hidden side effects; they are explicit events that can be inspected, routed, and repaired.

This is what makes the system codynamic rather than merely automated. The agents do not operate in isolation, each producing a local answer and handing it off as if the rest of the manuscript were irrelevant. Instead, each agent works against a shared work object and a shared runtime memory. The hypervisor, section agents, gardener, references agent, and diagram agent each own a distinct responsibility, but their work is coordinated through durable task queues and append-only messages. That coordination allows the manuscript to evolve as a connected whole rather than as a collection of disconnected files.

The outline is therefore both a plan and a control surface. It records the intended hierarchy of chapters and sections, but it also serves as a source of machine-readable dependencies and narrative expectations. Section agents use that structure to decide what to draft, what to revise, and what to defer. The gardener uses it to compare a section against its parent, siblings, and downstream uses. The references agent uses it to determine whether a claim needs support. The diagram agent uses it to decide whether a visual would materially improve the section. In each case, the outline is not merely descriptive; it participates in the system's behavior.

The same is true of artifacts. Imported source material, generated LaTeX payloads, compiled outputs, logs, and media files are not treated as disposable intermediates. They are durable evidence of the manuscript's state at a given moment. Because the system preserves these artifacts separately, it can reason about how a section was produced, what changed, and which downstream outputs were affected. That separation is essential to codynamic operation: the machine can revise one layer without losing sight of the others.

The thesis of this paper, then, is that long-form writing can be managed as a living technical system when the manuscript is represented as a coordinated set of evolving objects. The book machine works because it keeps those objects in sync: the outline informs the agents, the agents update the tasks, the tasks produce artifacts, the artifacts expose new dependencies, and the compiled outputs feed back into further revision. The result is not a one-time generation pipeline but a maintained authoring environment in which structure and content co-develop.

In the sections that follow, this paper treats the current repository as the concrete instance of that thesis. The repository is not just the place where the book lives; it is the operational evidence that the codynamic model is working.
